% ============================================================
% Chapter 17 — Research Methodology
% Source: PSYCON Engineering Design Document, Vol. 4, Ch. 17
% ============================================================
\section{Research Methodology}
\label{sec:research-methodology}

\subsection{Research Objective}
\label{sec:research-objective}

The objective of PSYCON is to investigate whether multimodal physiological and speech-derived features can be combined to estimate indicators relevant to psychological well-being in a research setting. The system is intended as a research and screening prototype, not a diagnostic medical device.

\subsection{Research Questions}
\label{sec:research-questions}

\textbf{Primary questions:}
\begin{itemize}
    \item Can physiological and speech features improve predictive performance compared with either modality alone?
    \item How robust are predictions under different environmental conditions?
    \item How does motion affect signal quality?
    \item What is the minimum recording duration needed for stable feature extraction?
\end{itemize}

\subsection{Study Design}
\label{sec:study-design}

\Cref{fig:study-workflow} shows the suggested end-to-end study workflow.

\begin{figure}[H]
\centering
\begin{tikzpicture}[
    font=\scriptsize,
    node distance=0mm,
    field/.style={draw, rounded corners=2pt, fill=blue!6, minimum width=3.4cm, minimum height=0.6cm, align=center, line width=0.5pt}
]
\node[field] (s1) {Participant Recruitment};
\node[field, below=5mm of s1] (s2) {Consent};
\node[field, below=5mm of s2] (s3) {Hardware Calibration};
\node[field, below=5mm of s3] (s4) {Data Collection};
\node[field, below=5mm of s4] (s5) {Data Quality Review};
\node[field, below=5mm of s5] (s6) {Feature Extraction};
\node[field, below=5mm of s6] (s7) {Model Training};
\node[field, below=5mm of s7] (s8) {Validation};
\node[field, below=5mm of s8] (s9) {Analysis};
\node[field, below=5mm of s9] (s10) {Reporting};

\foreach \i/\j in {s1/s2, s2/s3, s3/s4, s4/s5, s5/s6, s6/s7, s7/s8, s8/s9, s9/s10}
  \draw[-{Stealth[length=2mm]}, line width=0.6pt] (\i.south) -- (\j.north);
\end{tikzpicture}
\caption{Suggested end-to-end study workflow, from participant recruitment and consent through data collection, feature extraction, model training and validation, to final analysis and reporting.}
\label{fig:study-workflow}
\end{figure}

\subsection{Participants}
\label{sec:participants}

For each participant, the following should be recorded:
\begin{itemize}
    \item Anonymous participant ID
    \item Age group (if relevant)
    \item Biological sex (optional, if ethically approved)
    \item Recording date/time
    \item Session duration
    \item Environmental conditions
    \item Device firmware version
\end{itemize}

\begin{warningbox}
Avoid collecting personally identifying information unless necessary and approved.
\end{warningbox}

\subsection{Data Collection Protocol}
\label{sec:data-collection-protocol}

For every session:
\begin{enumerate}
    \item Verify hardware operation.
    \item Confirm battery level.
    \item Attach wrist module.
    \item Position audio module.
    \item Start synchronized recording.
    \item Monitor system status.
    \item End recording.
    \item Save data.
    \item Backup data.
    \item Review data quality.
\end{enumerate}

\subsection{Recorded Variables}
\label{sec:recorded-variables}

\Cref{tab:recorded-variables} lists the variables recorded for each session, organized by category.

\begin{table}[H]
\centering
\caption{Recorded variables by category.}
\label{tab:recorded-variables}
\small
\begin{tabularx}{\columnwidth}{@{}l Y@{}}
\toprule
\textbf{Category} & \textbf{Variables} \\
\midrule
Physiological & Heart rate; PPG waveform; GSR; temperature; motion; ambient light; battery status \\
Speech & Audio recording (if consented); extracted acoustic features; extracted language features \\
System & Timestamp; session ID; firmware version; sensor status; error logs \\
\bottomrule
\end{tabularx}
\end{table}

\subsection{Inclusion Criteria}
\label{sec:inclusion-criteria}

\textbf{Example criteria:}
\begin{itemize}
    \item Adult participants (or appropriate guardian consent for minors).
    \item Ability to provide informed consent.
    \item Ability to wear the device comfortably.
\end{itemize}

\subsection{Exclusion Criteria}
\label{sec:exclusion-criteria}

\textbf{Example criteria:}
\begin{itemize}
    \item Damaged skin at electrode location.
    \item Implanted medical devices, if required by the safety review.
    \item Participants unwilling to consent to physiological or audio recording.
\end{itemize}
